\section{The complete critical, cyclic-four, and carrier mixture}
\label{sec:carrier-mixture}
All groups and all direct-product coordinates in this section are
labelled. Write $\operatorname{SD}(G_1,\ldots,G_t)$ for the number of
subgroups full on every displayed coordinate, and put
\[
 \begin{aligned}
 \mathcal K_8&=\{8T18,8T26,8T27,8T28,8T29,8T31,8T35\},\\
 \mathcal K_{16}&=\{16T1082,16T1083,16T1084,16T1332,16T1547\}.
 \end{aligned}
\]
The names specify the literal permutation actions in the finite data.
A carrier of degree eight has weight one, and one of degree sixteen
has weight two. A carrier word of weight $T$ has degree $8T$.
Let $D$ be a critical word of elementary quotient rank $R$.

\subsection{Three simultaneous parameters on a full projection}
For a finite $2$-group $H$ set
\[
 d(H)=\dim H/\Phi(H),\qquad
 \rho(H)=\max_{M\lhd H,\,M\leq H'}d_H(M),\qquad
 d_H(M)=\log|M:M^2[M,H]|.
\]
The invariant $\tau(H)$ is the one in \eqref{eq:tau}; the five-term
sequence identifies it with $d_H(\Phi(H))$.

\begin{lemma}[Joint transition]\label{lem:carrier-transition}
Let $K\leq L\times G$ be full on both factors and let $N=K\cap G$.
Put
\[
 k=d_G(N),\quad n=\log|N|,\quad
 m=\max_{M\lhd G,\,M\leq N\cap G'}d_G(M),\quad
 R_N=N^2[N,G],\quad a_2=d_G(R_N).
\]
Then $\delta=d(K)-d(L)$ and
$\epsilon=\max(0,\rho(K)-\rho(L))$ satisfy
\begin{equation}\label{eq:carrier-polygon}
 0\leq\delta\leq k,\qquad 0\leq\epsilon\leq m,
 \qquad\delta+\epsilon\leq n,
 \qquad\tau(K)-\tau(L)\leq k-\delta+a_2.
\end{equation}
\end{lemma}
\begin{proof}
The map on Frattini quotients is onto, and
$\delta=\log|N:N\cap\Phi(K)|$. Fullness gives
$[N,K]=[N,G]$, so $R_N\leq N\cap\Phi(K)$ and $\delta\leq k$.
For $M\lhd K$ with $M\leq K'$, let $M_0$ be its projection to $L$
and $M_1=M\cap N$. Lifting normal generators gives
\[
 d_K(M)\leq d_L(M_0)+d_K(M_1).
\]
Here $M_1\lhd G$, $M_1\leq N\cap G'$, and
$d_K(M_1)=d_G(M_1)\leq m$. Also
$M_1\leq N\cap\Phi(K)$, so
$d_K(M_1)\leq n-\delta$. Maximization proves the two bounds on
$\epsilon$.

The map $\Phi(K)\to\Phi(L)$ is onto. After quotienting by squares
and ambient commutators, its kernel is generated by
$N\cap\Phi(K)$. Indeed products of squares and commutators in
$\Phi(L)$ lift to such products in $\Phi(K)$. Consequently
\[
 \tau(K)\leq\tau(L)+d_K(N\cap\Phi(K)).
\]
The quotient $(N\cap\Phi(K))/R_N$ is trivial elementary of dimension
$k-\delta$, and the image of $R_N$ has dimension at most $a_2$.
This proves the last inequality.
\end{proof}

For a row $r=(k,n,m,a_2,c,g)$ define
\[
 h_r(x,y)=\max\{x\delta+y\epsilon:
       0\leq\delta\leq k,\ 0\leq\epsilon\leq m,
       \delta+\epsilon\leq n\},\qquad q_r=\max(m,a_2).
\]
The quantity $a_2$ cannot in general be replaced by $m$.
The following criterion makes that replacement available for an entire
ambient group and reduces the radical certificates to generator tests.

\begin{lemma}[Finite radical certificates]\label{lem:carrier-finite-radical}
Let $G=\langle s_1,\ldots,s_u\rangle$ be a finite $2$-group and
$N=\langle x_1,\ldots,x_v\rangle\lhd G$. Then
\[
 R_N=\bigl\langle x_i^2,\ [x_i,s_j]:
       1\leq i\leq v,\ 1\leq j\leq u\bigr\rangle^G.
\]
If $\Phi(G)\leq G'$, then $R_N\leq N\cap G'$ and $a_2\leq m$
for every such $N$. If also $G'\leq N$, then $m=\rho(G)$ and
$G/N$ is abelian.
\end{lemma}
\begin{proof}
Write $U$ for the displayed normal closure. Its generators lie in
$R_N$, and $R_N\lhd G$, so $U\leq R_N$. In $G/U$, each $x_iU$
has square one and commutes with every $s_jU$, hence with the whole
quotient. The image of $N$ is therefore central elementary abelian,
which gives $N^2[N,G]\leq U$ and proves equality. An ordinary closure
can replace this normal closure when its normality in the original
$G$ has been certified.

Every square and ambient commutator from $N$ lies in $\Phi(G)$.
Under the stated containment this puts the same $G$-normal subgroup
$R_N$ inside $N\cap G'$. Its relative head is one of the terms in
the maximum defining $m$, so $a_2\leq m$. Finally $G'\leq N$ gives
$N\cap G'=G'$, proving the remaining assertions.
\end{proof}

\begin{lemma}[One centralizer controls all normal heads]
\label{lem:carrier-centralizer-head}
Let $G$ be a finite group, $p$ a prime, $N\lhd G$, and $u\in G$.
Write $R=N^p[N,G]$ and $d_{G,p}(N)=\log_p|N:R|$. Then
\[
 p^{d_{G,p}(N)}\leq |C_N(u)|.
\]
Consequently, for a finite $2$-group $G$ and any subgroup $S\leq G$,
\[
 2^{\max_{M\lhd G,\,M\leq S}d_G(M)}\leq |C_S(u)|,
 \qquad 2^{\rho(G)}\leq |C_{G'}(u)|.
\]
Neither $N$ nor $G'$ is required to be abelian.
\end{lemma}
\begin{proof}
The set map $f:N\to R$ given by $f(x)=[x,u]$ need not be a
homomorphism. If $f(x)=f(y)$, then $y^{-1}x\in C_N(u)$.
Choose one representative $y_r$ of each nonempty fibre. The map
\[
 x\longmapsto \bigl(f(x),y_{f(x)}^{-1}x\bigr)
\]
injects $N$ into $R\times C_N(u)$. Thus
$|N|\leq |R|\,|C_N(u)|$, proving the first bound after division
by $|R|$. Apply it to every original $G$-normal subgroup $M\leq S$
and use $C_M(u)\leq C_S(u)$. Taking the maximum gives the second
bound; putting $S=G'$ gives the last one.
\end{proof}

\begin{lemma}[Joint commutator radicals]\label{lem:carrier-joint-radicals}
Let $G$ be a finite $2$-group with $\Phi(G)=D=G'$, and put
\[
 V=G/D,\qquad E=D/(D^2[D,G]).
\]
For $\lambda\in E^*$ let $b_\lambda$ be the alternating form on $V$
given by $b_\lambda(xD,yD)=\lambda([x,y])$.
For every original normal subgroup $N\geq D$, write $W=N/D$ and
\[
 J_W=\{\lambda\in E^*: W\leq\operatorname{rad}(b_\lambda)\}.
\]
Then
\begin{equation}\label{eq:carrier-joint-head}
 d_G(N)\leq\dim W+\dim J_W.
\end{equation}
Let $q,a$ be nonnegative integers. If the forms associated with every
$q+1$ independent characters in
$E^*$ have common radical of dimension at most $a$, then
\[
 \dim W>a\quad\Longrightarrow\quad d_G(N)\leq\dim W+q.
\]
In particular, zero common radicals for independent character pairs give
$d_G(N)\leq\dim W+1$ for $W\ne0$. If each nonzero form has radical
dimension at most $a$, then $d_G(N)\leq\dim W$ whenever $\dim W>a$.
The subgroup $D$ need not be abelian.
\end{lemma}
\begin{proof}
Elements of $E^*$ are exactly the $G$-invariant characters of $D$.
The commutator identities and this invariance make
$\lambda([x,y])$ additive in both original variables. The resulting
maps vanish on $G'$ in each variable, so they descend to the displayed
alternating forms on $V$.

Restriction of a $G$-invariant character of $N$ to $D$ has kernel
the characters of $N/D$, of dimension $\dim W$. Every character in
the restriction image annihilates $[N,G]$, so that image lies in
$J_W$. Rank--nullity proves \eqref{eq:carrier-joint-head}. This uses
only a necessary condition on the restriction image; no assertion
that every member of $J_W$ extends is needed.

If $\dim J_W>q$, choose $q+1$ independent characters in $J_W$.
The same $W$ lies in all their radicals, contradicting
$\dim W>a$. Thus $\dim J_W\leq q$. The two special cases follow by
taking $(q,a)=(1,0)$ and $q=0$, respectively.
\end{proof}

\begin{corollary}[Star commutator family]\label{cor:carrier-star-family}
In the setting of Lemma~\ref{lem:carrier-joint-radicals}, suppose that
there are linear identifications $V=U\oplus\F_2$ and $E^*=U^*$ under
which the actual commutator forms are
\[
 b_\lambda((u,t),(v,s))=\lambda(u)s+t\lambda(v).
\]
Then every original normal subgroup $N\geq D$ satisfies
\[
 d_G(N)\leq\max\{\dim U,\dim(N/D)\}.
\]
\end{corollary}
\begin{proof}
Put $W=N/D$. If $W$ contains $(u,t)$ with $t\ne0$, then every
$\lambda\in J_W$ satisfies $t\lambda(v)=0$ for all $v\in U$, by
testing against $(v,0)$. Thus $J_W=0$.
Otherwise $W\leq U\oplus0$, and first projection identifies $W$
with a subspace $S\leq U$. Testing against $(0,1)$ shows that
$J_W=S^\perp\leq U^*$, so $\dim W+\dim J_W=\dim U$.
Both cases give the result through \eqref{eq:carrier-joint-head}.
The required identifications concern the complete original character
space and the actual quotient; ranks of selected numerical forms
alone do not supply them.
\end{proof}

\begin{lemma}[Derived-intersection directions]\label{lem:carrier-intersection-directions}
Keep $G,D,V,E$ as in Lemma~\ref{lem:carrier-joint-radicals}, but let
$N\lhd G$ be arbitrary. Put
\[
 B=N\cap D,\qquad W=ND/D,\qquad
 L_B=\{\lambda\in E^*: \lambda|_B=0\}.
\]
Then $L_B\leq J_W$ and $|N|=|B|2^{\dim W}$.
If $B<D$, then $L_B\ne0$. Consequently, a radical bound $a$ for
every nonzero $b_\lambda$ gives $\dim W\leq a$.
If independent pairs have zero common radical, then
$\dim L_B\geq2$ implies $N\leq D$.
For the star family of Corollary~\ref{cor:carrier-star-family},
$L_B\ne0$ gives the sharper bound
\[
 \dim W+\dim L_B\leq\dim U.
\]
\end{lemma}
\begin{proof}
Original normality gives $[N,G]\leq N\cap D=B$, so every character
in $L_B$ annihilates all commutators from $W$. The second isomorphism
theorem identifies $W$ with $N/B$, proving the order identity.
If $B<D$, choose a maximal proper $G$-normal subgroup $J$ of $D$
containing $B$. Then $D/J$ is a minimal nontrivial normal subgroup
of the finite $2$-group $G/J$, and hence has order two and is
central. Its nonzero character gives a member of $L_B$.

The single-form bound follows by choosing a nonzero member of $L_B$.
Two independent members force $W=0$ under the stated pair hypothesis.
For the star family, a nonzero member of $L_B$ forces every vector
of $W$ to have last coordinate zero. First projection identifies $W$
with a subspace of the common annihilator of $L_B$ in $U$, whose
dimension is $\dim U-\dim L_B$. This proves the last inequality.
These are constraints from characters of $D$ vanishing on $B$;
the extendibility and power constraints on characters of $B$ remain
separate.
\end{proof}

\begin{lemma}[Exact annihilator and center detection]
\label{lem:carrier-annihilator-center}
Keep $G,D,V,E$ as above and write $R=D^2[D,G]$.
For every subgroup $B\leq D$, the common kernel in $D$ of $L_B$
is $BR$. If $R\leq B$, put
\[
 C_B=\{x\in G:[x,G]\leq B\},\qquad
 U_B=\bigcap_{\lambda\in L_B}\operatorname{rad}(b_\lambda).
\]
Then $B\lhd G$, $D\leq C_B$, and the original quotient map induces
\[
 C_B/D\cong U_B,\qquad |C_B|=|D|2^{\dim U_B}.
\]
Suppose in addition that independent character pairs have zero common
radical and that $N\lhd G$ satisfies neither $N\leq D$ nor $D\leq N$.
For $B=N\cap D$, the space $L_B$ has a unique nonzero member
$\lambda$. If $R\leq B$, then, with
$w=\dim(ND/D)$ and $r=\dim\operatorname{rad}(b_\lambda)$,
\begin{equation}\label{eq:carrier-crossing-center}
 |D:B|=2,\qquad 1\leq w\leq r,\qquad
 |Z(G/N)|=2^{r+1-w},\qquad |(G/N)'|=2.
\end{equation}
\end{lemma}
\begin{proof}
The invariant characters of $D$ are the complete dual of $D/R$.
The ordinary double-annihilator identity in this vector space says
that their common kernel on $B$ is the inverse image of $BR/R$,
which is $BR$. This does not identify $B$ itself unless $R\leq B$.

Under that containment, $B/R$ lies in the central elementary subgroup
$D/R$ of $G/R$, so $B$ is normal in $G$.
For every original $x,y\in G$, detection gives
$[x,y]\in B$ exactly when $\lambda([x,y])=0$ for every
$\lambda\in L_B$. Thus $x\in C_B$ exactly when $xD\in U_B$.
In particular $D\leq C_B$, and the asserted equivalence and order
formula follow from the original quotient map.

In the crossing case, Lemma~\ref{lem:carrier-intersection-directions}
gives $L_B\ne0$ and $W=ND/D\ne0$. Two independent members of $L_B$
would force $W=0$, so $\dim L_B=1$. Its unique nonzero member
detects $BR$; when $R\leq B$, this is its literal kernel $B$ in $D$,
of index two. Also $W\leq\operatorname{rad}(b_\lambda)$, so
$1\leq w\leq r$. Since $|N|=|B|2^w$ and
$Z(G/N)=C_B/N$, the center formula follows. Finally the original
derived group maps onto $(G/N)'$ with kernel $B$, proving its order.
\end{proof}

\begin{lemma}[Saturation and exact crossing heads]
\label{lem:carrier-saturation-crossing}
Keep the notation of Lemma~\ref{lem:carrier-annihilator-center}.
Suppose $\dim E\geq2$, independent character pairs have zero common
radical, and every original $G$-normal subgroup $M\leq D$ satisfies
\[
 M\leq R\quad\hbox{or}\quad R\leq[M,G].
\]
Then every crossing normal $N$ satisfies, with $B=N\cap D$,
\[
 [N,G]=B=N^2[N,G],\qquad d_G(N)=\dim(ND/D).
\]
For every original normal $B\leq D$ with $B\not\leq R$,
\[
 B^2[B,G]=R,\qquad
 \max_{M\lhd G,\,M\leq B}d_G(M)
 =\max\left\{\max_{M\lhd G,\,M\leq R}d_G(M),d_G(B)\right\},
\]
where $d_G(M)=\dim\operatorname{Hom}(M,\F_2)^G$.
\end{lemma}
\begin{proof}
Put $H=[N,G]\leq B$. If $H\leq R$, every character of $D$
annihilates $H$, so every form annihilates $W=ND/D$ against $V$.
Two independent characters would then force $W=0$, a contradiction.
Thus saturation gives $R\leq H$. The full space $L_H$ is the joint
annihilator of $W$, and the pair hypothesis makes its dimension at most
one. Since $0\ne L_B\leq L_H$, these spaces are equal. Exact detection
now gives $H=B$. Every original square in $N$ lies in $N\cap D$,
since $G/D$ has exponent two, proving the relative radical identity
and the head formula.

If $M\leq D$ and $M\not\leq R$, saturation and monotonicity of the
relative radical give $R\leq[M,G]\leq M^2[M,G]\leq R$.
For $M\leq B$ outside $R$, the resulting quotient inclusion
$M/R\leq B/R$ gives $d_G(M)\leq d_G(B)$. Subgroups inside $R$
contribute the other maximum. Both terms are attained within $B$,
which proves equality without enumerating its normal subgroups.
\end{proof}

For the original pair-family masters $16T1082$, $16T1083$, $16T1084$
and $16T1547$, the saturation hypothesis is certified by original
commutator words on their $16$ or $64$ derived rows. Every word has
positive commutator depth, so an element of a normal subgroup outside
$R$ generates all of $R$ through its original ambient commutators.
The same finite certificates give saturation for $16T1332$, whose
star family is treated separately from the pair hypothesis.
For the four pair-family masters, every nonzero form has radical
dimension two, $d_G(D)=3$, and
$\max_{M\lhd G,\,M\leq R}d_G(M)\leq2$.
Consequently every crossing normal has exactly one of the following
rows, in the order $(k,n,m,a_2,c,g)$:
\[
\begin{array}{c|cc}
 G & w=1 & w=2\\ \hline
 16T1082,\ 16T1083,\ 16T1084
 &(1,4,2,2,2,1)&(2,5,2,2,1,1)\\
 16T1547 &(1,6,2,2,2,1)&(2,7,2,2,1,1).
\end{array}
\]
Indeed $B$ has index two in $D$, its relative radical is $R$, and
$d_G(B)=2$; hence $m=a_2=2$. The remaining coordinates follow from
Lemma~\ref{lem:carrier-annihilator-center} and $|N|=|B|2^w$.
These are exhaustive possibilities for crossing normals; no realization
of every listed row or identification of different normals is required.

For $16T1547$, the remaining normals inside $R$ also admit a uniform
argument. Here $|R|=8$, and the original commutator tests on its eight
rows give a central subgroup $Z$ of order two with
$R^2[R,G]=[R,G]=Z$. Every original normal $M\leq R$ either lies in
$Z$ or contains $Z$; if $M\not\leq Z$, the same saturation words give
$M^2[M,G]=Z$. Thus $M$ is $1$, $Z$, an order-four subgroup above $Z$,
or $R$. This argument requires no list or count of the order-four
subgroups. Their exact rows are
\[
\begin{array}{c|c}
 M &(k,n,m,a_2,c,g)\\ \hline
 1 &(0,0,0,0,1,6)\\
 Z &(1,1,1,0,2,5)\\
 Z<M<R &(1,2,1,1,1,4)\\
 R &(2,3,2,1,3,3).
\end{array}
\]
For completeness, the original form separation first gives $C_M\leq D$.
If $M<R$, saturation strengthens this to $C_M\leq R$.
The identity $[R,G]=Z$ then gives $C_M=R$ for $Z\leq M<R$;
the eight-row fixed tests give $C_1=Z$, and $C_R=D$.
These are the center orders in the table. The derived orders are
$|D|/|M|$, and the maximum defining $m$ ranges over every original
normal below $M$, using the same relative-radical identities.

\begin{lemma}[Crossing capacities for the full star family]
\label{lem:carrier-star-crossing}
Keep $G,D,V,E$ as above, and suppose the complete original form family
is the star family on $V=U\oplus\F_2$, with $E^*=U^*$ and $d=\dim U$.
Assume every original normal subgroup of $D$ is comparable with $R$.
For a crossing normal $N$, put $B=N\cap D$, $H=[N,G]$,
$w=\dim(ND/D)$ and $g=\dim L_B$. Then
\[
 R\leq H\leq B,\qquad d_G(N)+g\leq d,\qquad
 |(G/N)'|=2^g,\qquad |Z(G/N)|=2^{d-w}.
\]
\end{lemma}
\begin{proof}
The complete space $L_H$ is exactly the joint annihilator of $W=ND/D$.
Since it contains the nonzero space $L_B$, the star formula forces
$W\leq U\oplus0$. Its first projection identifies $L_H$ with the
annihilator of a $w$-dimensional subspace of $U$, so
$\dim L_H=d-w$. If $H\leq R$, then $L_H=E^*$, forcing $w=0$;
comparability therefore gives $R\leq H$.
Exact character detection and the order identity now give
\[
 |D|=2^{d-w}|H|=2^g|B|,\qquad |N|=2^w|B|.
\]
All original powers remain in the relative radical $N^2[N,G]$,
which contains $H$. Hence
$2^{d_G(N)}|H|\leq|N|$, and the displayed identities yield
$d_G(N)+g\leq d$.

The full common radical of $L_B\ne0$ is its annihilator in
$U\oplus0$, of dimension $d-g$. Thus
$|C_B|=|D|2^{d-g}$ and $|Z(G/N)|=|C_B|/|N|=2^{d-w}$.
The image of $D$ in $G/N$ has order $|D:B|=2^g$.
This proof uses neither a one-dimensional parameter space nor an equality
between $H$ and the full relative radical of $N$.
\end{proof}

\begin{lemma}[The remaining original normal branches]
\label{lem:carrier-normal-branches}
For each of the five degree-sixteen masters put $D=G'$,
$R=D^2[D,G]$, $d=\log|D|$, $s=\log|G/D|$ and $e=d_G(D)$.
Their parameters are
\[
\begin{array}{c|rrrr}
G&d&s&e&\log|R|\\\hline
16T1082,\ 16T1083,\ 16T1084&4&6&3&1\\
16T1332&6&5&4&2\\
16T1547&6&6&3&3.
\end{array}
\]
If $D\leq N$, with $w=\dim(N/D)$, then $0\leq w\leq s$ and
the six invariants are bounded coordinatewise by
\[
 (\max(e,w),d+w,e,e,s-w,0).
\]
Here the order, the maximum $m=e$, and both quotient slopes are exact.
For $R<N<D$ the exhaustive exact alternatives are
\[
\begin{array}{c|l}
G&(k,n,m,a_2,c,g)\\\hline
16T1082,\ 16T1083,\ 16T1084
 &(1,2,1,1,2,2),\ (2,3,2,1,3,1)\\
16T1547 &(1,4,2,2,2,2),\ (2,5,2,2,3,1)\\
16T1332 &(t,2+t,t,1,4,4-t),\quad t\in\{1,2,3\}.
\end{array}
\]
Inside $R$, the first three masters have the two exact rows
$(0,0,0,0,1,4)$ and $(1,1,1,0,3,3)$.
For $16T1332$ there is a central subgroup $Z$ of order two, and every
original normal inside $R$ is $1,Z$ or $R$, with respective rows
\[
 (0,0,0,0,1,6),\qquad(1,1,1,0,1,5),\qquad(1,2,1,1,4,4).
\]
Together with the four small-radical rows of $16T1547$ and the crossing
results above, these alternatives cover every original normal subgroup.
\end{lemma}
\begin{proof}
For $D\leq N$, the image dimension gives $|N|=2^{d+w}$ and the
abelian quotient has order $2^{s-w}$. The original derived-normal
maximum is $e$: its upper bound follows from the checked original
centralizer ranks, and $D$ itself attains the lower bound. The full-form
head estimate gives $d_G(N)\leq\max(e,w)$, while
$N^2[N,G]\leq N\cap D$ gives $a_2\leq m=e$.

For $R<N<D$, nonempty original commutator saturation gives
$N^2[N,G]=R$. Put $t=d_G(N)$ and $\ell=\dim L_N$.
Then $t,\ell>0$, $t+\ell=e$, $|N|=|R|2^t$, and
$m=\max(m(R),t)$. In a pair family, $e=3$; the full common radical
of $L_N$ has dimension zero when $\ell=2$, and dimension two when
$\ell=1$. Lemma~\ref{lem:carrier-annihilator-center} gives
$c=\ell+\dim U_N$ and $g=\ell$.
In the complete star family the nonzero full parameter space $L_N$
has common radical dimension $4-\ell=t$, so $c=4$ and $g=4-t$.
The excluded case $\ell=0$ is $N=D$ and belongs to the first branch.

For the first three masters, $|R|=2$ makes $R$ central, its relative
radical trivial, and its head and normal-head maximum both one.
Full pair-form separation and saturation give $C_1=R$ and $C_R=D$,
proving the two small rows. For $16T1332$, the four original radical
rows, with all five original generators acting, give
$R^2[R,G]=[R,G]=Z$, with $|R|=4$ and $|Z|=2$.
Every original normal $M\leq R$ lies in $Z$ or contains $Z$;
outside $Z$ it has relative radical $Z$ and head at most one.
Cardinality therefore forces $M=R$. The same tests and full star
separation give $C_1=Z$, $C_Z=R$, and $C_R=D$, yielding the three
rows and $m(R)=d_G(R)=1$.

Finally split an arbitrary $N\lhd G$ first by $D\leq N$, then by
$N\leq D$. In the latter case saturation makes $N$ comparable with
$R$, giving either $N\leq R$ or $R<N<D$. The remaining case is
crossing. This proof never counts normals by counting their possible
rows and does not assume that every displayed possibility is realized.
\end{proof}

There is a useful simplification when comparing these rows numerically.
The feasible pairs satisfy $\delta\leq k$, $\epsilon\leq m$ and
$\delta+\epsilon\leq n$, so their polygon depends on $n$ only through
$\min(n,k+m)$. Thus a comparison row $(K,N_0,M,A,C,G_0)$ dominates
all pair supports and both marks whenever
\[
 k\leq K,\quad m\leq M,\quad a_2\leq A,\quad
 \min(n,k+m)\leq N_0,\quad c\leq C,\quad g\leq G_0.
\]
This comparison preserves the original order $n$ in the group-theoretic
identities. It simply includes its entire feasible polygon in the
comparison polygon. In particular it remains valid for a negative
first mark, with a nonnegative terminal multiplier.

For $Q=G/N$ take $(c,g)=(\log|Q|,0)$ if $Q$ is abelian, and
$(c,g)=(\log|Z(Q)|,\log|Q'|)$ otherwise. For a full subdirect
product $L$ of the bounded carrier factors,
\begin{equation}\label{eq:carrier-epi}
 |\Epi(L,Q)|\leq C(d(L)+2)^C2^{c d(L)+g\rho(L)}.
\end{equation}
The following elementary bound proves \eqref{eq:carrier-epi} for every
finite $2$-group source, without a restriction on its direct-product
factors. Its constant is absolute on the fixed target alphabet.

\begin{proposition}[Structured epimorphisms]\label{prop:structured-epi}
Let $G,Q$ be finite $2$-groups, and put $d=d(G)$,
$\rho=\rho(G)$ and $q=\log|Q|$. Then
\begin{equation}\label{eq:structured-epi-explicit}
 |\Epi(G,Q)|\leq
 (d+1)^q\,|Q'|^{\rho}\,|Q|^{|Q'|^q}\,|Z(Q)|^d.
\end{equation}
In particular, \eqref{eq:carrier-epi} holds with
\[
 C=C_Q=\max\{q,|Q|^{|Q'|^q}\}\geq1,\qquad
 c=\log|Z(Q)|,\qquad g=\log|Q'|.
\]
These choices also give the stated abelian case.
\end{proposition}
\begin{proof}
Choose an actual generating tuple of $G$ of length $d$, by lifting a
basis of $G/\Phi(G)$. For every epimorphism to $Q$, at most $q$ of
these same generators have images generating $Q$: choose a basis
from their images in $Q/\Phi(Q)$ and lift back to the selected
indices. Padding by identities gives at most $(d+1)^q$ ordered
tuples. Fix one such tuple $x=(x_1,\ldots,x_q)$ and count the maps
$f$ for which $f(x)$ generates $Q$.

Set $D=[\langle x\rangle^G,G]\leq G'$ and
$H=D\cap\ker f$. These are normal subgroups of the original $G$,
and $f(D)=Q'$, so $[D:H]=|Q'|$. There are at most $|Q'|^{\rho}$
possible $H$. Indeed, for every $M\lhd G$ inside $G'$, the number
of $G$-normal subgroups of index at most $2^a$ in $M$ is at most
$2^{a\rho}$. To see this, every proper normal interval in a finite
$2$-group admits a bottom step of index two: choose a central
involution in its nontrivial normal quotient. Each normal $J\leq M$
has at most $2^{d_G(J)}-1\leq2^{\rho}-1$ normal index-two children,
since each is the kernel of a nonzero $G$-invariant character of $J$.
Keeping the previous subgroups as well as these children multiplies
the cumulative count by at most $2^{\rho}$ at each step, starting
with the single subgroup $M$.

Now fix $H$ and let $C$ be the inverse image in $G$ of the
centralizer of the tuple $xH$ in $G/H$. The conjugates of that tuple
lie coordinatewise in $x_iD/H$, so
\[
 [G:C]\leq|D/H|^q=|Q'|^q.
\]
Every counted $f$ sends $C$ into $Z(Q)$, because it kills $H$ and
$f(x)$ generates $Q$. Thus $G\to Q/Z(Q)$ is constant on left
$C$-cosets, giving at most $|Q|^{|Q'|^q}$ quotient maps. Normality
of $C$ is not required. In a fixed quotient-map fibre, choose one
map $f_0$ when the fibre is nonempty. The pointwise differences
$f(g)f_0(g)^{-1}$ are central and form a homomorphism $G\to Z(Q)$;
they determine $f$. There are at most $|Z(Q)|^d$ such
homomorphisms, by the original generating tuple. Empty fibres
contribute zero. Multiplying the tuple, normal-subgroup,
quotient-map and central-fibre bounds proves
\eqref{eq:structured-epi-explicit}. Finally
$(d+1)^q\leq(d+2)^{C_Q}$ and $|Q|^{|Q'|^q}\leq C_Q$.
\end{proof}

\subsection{The finite rows and the recurrence}
Use the masters $X=8T26$, $J=8T27$, $P=8T35$ and the five degree-sixteen
carriers. The finite normal-kernel certificate supplies their actual
rows and their multiplicities. It includes every literal normal kernel,
by the central-involution closure argument in
Lemma~\ref{lem:binary-coverage} below. The numerical reduction stated
and certified in Appendix~\ref{sec:carrier-cones} gives the matrix
\begin{equation}\label{eq:carrier-B}
 B=\frac1{64}\begin{pmatrix}
 5&6&34&28&32&30\\6&0&44&36&40&36\\
 34&44&0&16&20&40\\28&36&16&12&16&32\\
 32&40&20&16&16&32\\30&36&40&32&32&24
 \end{pmatrix}.
\end{equation}
\[
 \alpha=(5/8,3/4,0,1/4,1/4,1/2),\quad
 \beta=(3/8,1/2,0,1/4,1/4,1/2).
\]
Here rows have physical scale $s_r=1$ or $2$. Before domination their
pair coefficient is
\[
 \frac{\max\{h_r(c_s,g_s),h_s(c_r,g_r)\}}{s_rs_s},
\]
and their marked coefficients are $k_r/s_r,q_r/s_r$. Each dominated row is moved to a
row dominating all these coefficients simultaneously. Finally scale
its accumulated mass by four, obtaining $p_i\geq0$ with
$\sum_i p_i=4T$; this accounts for the factor $1/64$ in
\eqref{eq:carrier-B}.

The crossing branch of $16T1332$ can be assigned directly to the
sixth colour, without distinguishing its six possible pairs $(g,w)$.
Indeed Lemma~\ref{lem:carrier-star-crossing} and saturation give the
coordinatewise upper row
\[
 (4-g,\,6-g+w,\,\max(2,4-g),\,\max(2,4-g),\,4-w,\,g),
 \qquad g,w\geq1,\quad g+w\leq4.
\]
All these rows lie below $r_*=(3,8,3,3,3,3)$ at scale two.
For any displayed row $s=(k,n,m,a_2,c,g)$, with nonnegative slopes,
\[
 \max\{h_s(3,3),h_{r_*}(c,g)\}
 \leq\max\{3\min(n,k+m),\,3(c+g)\}.
\]
Checking this last expression on the forty-one displayed rows of
Appendix~\ref{sec:carrier-cones}, with their stated colours $j$ and
scales $s_s$, gives the bound $8(4s_s)B_{6j}$ in every column.
The self support is at most $18\leq64B_{66}=24$, and both marks
are $3\leq8\alpha_6=8\beta_6=4$.
This handles all crossing normals of that master at once, while
retaining their separate original weights. The remaining branches are
covered by Lemma~\ref{lem:carrier-normal-branches}. Its effective-polygon
comparison sends every pair-family row, every above-derived row, and
the small-radical rows into the twelve displayed H16 rows. The first
two star intermediate rows go to H16 rows $2$ and $4$.
The last star intermediate row is $r_\dagger=(3,5,3,1,4,1)$, at scale
two, and can also be assigned directly to colour six. Indeed
\[
 \max\{h_s(4,1),h_{r_\dagger}(c,g)\}
 \leq\max\{4k+m,\,3(c+g)\}.
\]
The forty-one displayed comparisons satisfy the same bound
$8(4s_s)B_{6j}$. Its self support is at most $15\leq24$ and its
pair support with $r_*$ is at most $18\leq24$; both marks are at most
$4$. Thus these two extra envelopes use the existing sixth colour
and do not change the cone matrix or its physical normalization.
No step replaces the original weighted normal family by a sum over
distinct numerical labels.

\begin{lemma}[The complete marked recurrence]\label{lem:carrier-recurrence}
For every such master word, the logarithm of its full-subdirect count
with $D$ is at most the maximum, over $0\leq j\leq R$,
$0\leq\ell\leq R/2$ and its row masses, of
\begin{equation}\label{eq:carrier-E}
 \mathcal E=\ell(R/2-\ell)+(j-\ell)(R-j)
   +\tfrac12p^{\mathsf T}Bp
   +\sum_i p_i\{\ell(\alpha_i+\beta_i)
                       +\alpha_i\max(j-2\ell,0)\},
\end{equation}
plus $O((R+T)\log(R+T+2))$.
\end{lemma}
\begin{proof}
Project first onto the complete carrier subgroup $H$. For this same
$H$, Proposition~\ref{prop:terminal-bound} gives the sum of the marks
$2^{(j-\ell)d(H)+\ell\tau(H)}$, with coefficient
\[
 \phi^{-2}\gauss c\ell2\,72^\ell2^{(j-\ell)(R-j)},
 \qquad c\leq R/2.
\]
For one carrier peel the lower projection $L$, literal axis $N$ and
literal epimorphism $L\to G/N$ recover the subgroup uniquely. Thus
\eqref{eq:carrier-epi} is used once, with no automorphism divisor.
Lemma~\ref{lem:carrier-transition} shows that an arbitrary current
mark $2^{x d+y\rho+z\tau}$, for $y,z\geq0$, incurs at most
\[
 z(k+q_r)+h_r(x-z,y).
\]
This remains true for negative $x$. For the initial mark
$(x,y,z)=(j-\ell,0,\ell)$ it is at most
$\ell(k+q_r)+k\max(j-2\ell,0)$.
Every later epimorphism contributes the nonnegative tilt $(c_s,g_s)$.
Subadditivity of the support function bounds the additional cost of
an unordered pair of peels by the maximum in the definition of the
pair coefficient. Summing pairs gives at most the quadratic term
$\tfrac12p^{\mathsf T}Bp$; adding diagonal self-pairs only enlarges it.
The simultaneous row domination gives the last term of
\eqref{eq:carrier-E}. Finally use
$\gauss c\ell2\leq\phi^{-1}2^{\ell(c-\ell)}$.

There are a fixed number of normal choices per coordinate, hence at
most $2^{O(T)}$ histories. The polynomial factors in
\eqref{eq:carrier-epi} cost $O(T\log(T+2))$; $72^\ell$, the Gaussian
constants and the outer $(j,\ell)$ sum cost $O(R+\log(R+2))$.
This proves the stated uniform error. The initial projection $H$ is
recovered intrinsically, so the argument attaches its critical factors
exactly once.
\end{proof}

\begin{theorem}[Critical--carrier reserve]\label{thm:carrier-cross}
Uniformly over all original carrier colour words of weight $T\geq1$,
\begin{equation}\label{eq:carrier-cross}
 \log\operatorname{SD}(D,\boldsymbol G)
 \leq\frac{(R+4T)^2}{4}-\frac{25}{82}RT-\frac{29}{164}T^2
           +O((R+T)\log(R+T+2)).
\end{equation}
The sum of their original action weights is
\begin{equation}\label{eq:carrier-colours}
 A_T=[z^T]\exp\left(\frac3{64}z+\frac{65}{688128}z^2\right).
\end{equation}
\end{theorem}
\begin{proof}
Appendix~\ref{sec:carrier-cones} proves, for every point of the two
exhaustive cones $j\geq2\ell$ and $j\leq2\ell$,
\[
 \mathcal E\leq R^2/4+\frac{139}{328}R\sum_i p_i
                    +\frac{627}{2624}\left(\sum_i p_i\right)^2.
\]
With $\sum_i p_i=4T$, this is \eqref{eq:carrier-cross} for masters.
There are fixed epimorphisms
\[
 X\to8T26,\qquad J\to8T27,8T28,\qquad
 P\to8T18,8T29,8T31,8T35.
\]
The inverse image under their product, acting identically on $D$ and
all degree-sixteen factors, injects each original full-subdirect
family into its master family. Its image recovers the original
subgroup, and fullness of each replaced coordinate follows from
surjectivity. This proves the bound for each original word.

The degree-eight normalizers have orders $192$ for $8T18$, $384$
for $8T31$, and $128$ for the other five; their reciprocals sum to
$3/64$. The degree-sixteen orders, in the displayed order, are
$344064,49152,49152,24576,98304$, whose reciprocals sum to
$65/688128$. Summing $\prod_U1/(a(U)^{m_U}m_U!)$ proves
\eqref{eq:carrier-colours}. Each injection is used profile by profile
with its original weight. Overlap between the images of different
profiles causes no identification of the original colours.
\end{proof}

\subsection{Adding every cyclic-four orbit}
Let $a$ be the number of regular $C_4$ orbits, let $T$ be the carrier
weight and put $N=R+2a+4T$. The exact original profile weight is
\begin{equation}\label{eq:mixture-weight}
 \frac{c_R}{8^a a!}A_T,
 \qquad c_R=[y^R]\exp(y/2+y^2/6+y^4/384).
\end{equation}

\begin{lemma}[Cyclic-four comparison and Hall bound]\label{lem:mixture-Hall}
For any finite $2$-group $B$ of order $2^u$,
\begin{align}\label{eq:C4-comparison}
 \operatorname{SD}(D,C_4^a,B)
 &\leq\sum_{k=0}^{2a}\binom{2a}{k}
                 \operatorname{SD}(D,C_2^k,B),\\
 \log|\operatorname{Sub}(C_2^R\times C_4^a\times B)|
 &\leq\frac{u^2+(R+a+u)^2+a^2}{4}
       +O((a+u)\log(a+u+2)+\log(R+2)).\label{eq:Hall-mixture}
\end{align}
In the first inequality $B$ may be a displayed word, with fullness
required on each of its factors.
\end{lemma}
\begin{proof}
For a central extension by $C_2$, the subgroups above a fixed subgroup
of the quotient consist of its unique full preimage and its complements
to $C_2$. The latter are absent or form a torsor of size
$|\Hom(H,C_2)|$. The split extension realizes the upper bound exactly.
Apply this successively to the central involutions in the $C_4$
coordinates, retaining the projection to $D\times B$ throughout.
After dropping fullness only on these coordinates, the comparison
group is $C_2^{2a}$. Its uniquely determined set of nonzero coordinate
projections has size $k$, giving \eqref{eq:C4-comparison}.
For the terminal upper bound, this support decomposition is unnecessary:
the terminal family already permits every image in its binary abelian
coordinates. Regroup the new $2a$ binary coordinates with the original
binary abelian part, retaining fullness on each nonabelian critical
coordinate and each carrier factor. The resulting terminal family has
critical rank $R+2a$, with no binomial prefactor. This regrouping concerns
the comparison family; the original orbit profile and its weight remain
unchanged.

The same central comparison gives, for every finite group $E$,
\[
 |\operatorname{Sub}(E\times B)|
 \leq |\operatorname{Sub}(E\times C_2^u)|.
\]
Indeed, choose a central subgroup $C\leq B$ of order two and compare
$E\times B\to E\times(B/C)$ with its split extension. Apply induction
to $B/C$, keeping $C_2\times E$ as the unchanged exterior, and regroup
the binary factors. This induction preserves every condition on the
exact $E$-image; it does not require $E$ to be abelian.

Set $E=C_4^a\times C_2^R$ and $r=R+u$. Reduce the $C_4$ coordinates
of $C_4^a\times C_2^r$ modulo two. Its kernel has binary dimension
$r+a$. Above an actual image $U$ of dimension $j$, the squares span
exactly the copy of $U$ in the doubled $C_4$ coordinates. The square
obstruction therefore has dimension $j$. The exact central-lift formula,
with $k$ the retained annihilator dimension, gives
\[
 |\operatorname{Sub}(C_4^a\times C_2^r)|
 =\sum_{j=0}^a\gauss aj2
    \sum_{k=0}^{r+a-j}\gauss{r+a-j}k2\,2^{jk}.
\]
Use $\gauss nk2\leq\phi^{-1}2^{k(n-k)}$, where
$\phi=\prod_{i\geq1}(1-2^{-i})$. The exponent of each summand is at most
\[
 j(a-j)+k(r+a-j-k)+jk
 =j(a-j)+k(r+a-k)\leq\frac{a^2+(r+a)^2}{4}.
\]
There are at most $(a+1)(r+a+1)$ summands. In particular we have the
stronger explicit bound
\begin{equation}\label{eq:Hall-mixture-explicit}
 |\operatorname{Sub}(C_2^R\times C_4^a\times B)|
 \leq \frac{(a+1)(R+a+u+1)}{\phi^2}
       2^{((R+a+u)^2+a^2)/4}.
\end{equation}
The logarithm of the prefactor is
$O((a+u)\log(a+u+2)+\log(R+2))$, uniformly also when $a=u=0$.
Adding the nonnegative $u^2/4$ recovers \eqref{eq:Hall-mixture}
with the stated constants.
\end{proof}

\begin{theorem}[Complete mixture]\label{thm:binary-mixture}
The normalized mass, relative to $L_{2N}$, of all fixed-point-free
binary subgroups whose orbit actions belong to the critical actions,
regular $C_4$, and the twelve carriers, with some noncritical orbit,
is $2^{-\Omega(N)}$. At a fixed profile the following stronger
bounds hold uniformly:
\begin{align}\label{eq:proved-mixture-envelopes}
 T\leq a/140,\ a\geq1:&\quad
 2^{-aN/25+O(a\log(N+2))},\notag\\
 T>a/140:&\quad
 2^{-(29/656)TN+O(N\log(N+2))}.
\end{align}
\end{theorem}
\begin{proof}
The carrier product $B$ has $u=\log|B|\leq7T$. Use the
exterior-preserving central comparison in Lemma~\ref{lem:mixture-Hall},
and regroup its unrestricted binary abelian coordinates as above.
The terminal reserve from Theorem~\ref{thm:carrier-cross}, applied
directly at critical rank $R+2a$, gives
\[
 \begin{aligned}
 \log\operatorname{SD}(D,C_4^a,B)
 &\leq N^2/4-\frac{25}{82}(R+2a)T-\frac{29}{164}T^2
                  +O(N\log(N+2))\\
 &\leq N^2/4-\frac{29}{656}TN+O(N\log(N+2)).
 \end{aligned}
\]
The original noncritical profile weights have a uniform finite-sum
bound without evaluating their normalizers. If $w_i$ is the original
normalizer order of colour $i$, then $w_i\geq1$, and any finite set
$\mathcal M$ of multiplicity profiles satisfies
\[
 \sum_{\mathbf m\in\mathcal M}
    \prod_{i=1}^{13}\frac1{w_i^{m_i}m_i!}
 \leq\prod_{i=1}^{13}\exp(1/w_i)\leq e^{13}.
\]
All thirteen colours remain separate, including those with the same
master. Together with $c_R/c_N\leq(2N)^{2a+4T}$, division by $c_NG_N$
and insertion of \eqref{eq:mixture-weight} proves the second envelope.

For the other envelope first project onto an arbitrary actual
$H\leq C_4^a\times B$. Its sectional binary rank is at most $a+u$:
intersect an elementary section with $H\cap C_4^a$ and project the
quotient to $B$. Hence $\tau(H)\leq a+u$.
Proposition~\ref{prop:terminal-bound}, with the nonnegative $\delta$
and $d/2$ dropped, gives
\[
 N_D(H)\leq X_R(d(H))
 2^{h_R(a+u+\log72)+O(\log(R+2))},\qquad
 h_R(x)=\max_{0\leq\ell\leq R/2}(x\ell-3\ell^2/4)\leq x^2/3.
\]
The graph identity
$\sum_{H\leq C_4^a\times B}X_R(d(H))
=|\operatorname{Sub}(C_2^R\times C_4^a\times B)|$
then allows a single use of \eqref{eq:Hall-mixture}.
Put $N_0=R+2a$. Subtracting the resulting quadratic exponent from
$N_0^2/4$ leaves exactly
\[
 aR/2+a^2/6-uR/2-7au/6-5u^2/6.
\]
If $T\leq a/140$, then $u\leq a/20$, so this is at least
$(19/40)aR+(17/160)a^2\geq aN_0/20\geq aN/25$;
here $N\leq71N_0/70$. The fixed $\log72$ correction costs $O(a+1)$,
and all profile and coefficient costs are $O(a\log(N+2))$.
This proves the first envelope, including $T=0$.

Let $C=2a+4T$ be the noncritical half-support. For this finite alphabet,
the entire range $0<C\leq\sqrt N/4$ has a direct estimate. Each original
carrier action has binary character rank at most three times its physical
scale: apply the certified head bound to the top normal subgroup and
inflate characters along its fixed master epimorphism. This is strictly
less than half its original degree. The regular $C_4$ has rank one.
For an actual full noncritical tail $J$, restrict to one of its original
blocks. The restriction kernel acts faithfully on the original complement,
so the half-degree bound for that kernel preserves the local gap and gives
$d_2(J)\leq C-1$, even when $J$ is a proper subdirect product. Applied to
the actual evaluation kernel, the same bound gives $\tau(J)\leq C$.
These marks concern $J$ itself, before any master pullback.

To count such tails, the master pullback is injective, and its source
has order $2^q$ with $q=2a+u\leq7C/4$. The central comparison above
bounds its subgroup count by $G_q$. The terminal Hall estimate therefore
bounds the model count by
\[
 2\phi^{-3}(R+1)(c+1)\,X_R(C-1)\,G_q\,
       2^{(C+7)^2/3},
\]
where $c$ is the number of nonabelian critical factors. The elementary
Gaussian estimates give
$X_R(C-1)/G_N\leq\phi^{-2}(R+1)2^{-N/2+1/2}$ and
$G_q\leq\phi^{-1}(q+1)2^{q^2/4}$. Sum the critical profiles at their
exact rank $R=N-C$, retaining $c_R$, and the noncritical profiles with
their original weights as above. The logarithm of the normalized mass
at this support is at most
\[
 -\frac N2+\frac12+\frac{211}{192}C^2+\frac{14}{3}C+\frac{49}{3}
   +C\log(2N)+O(\log(N+2)).
\]
For $C\leq\sqrt N/4$ the quadratic terms leave at most
$-1325N/3072$, and the remaining error is uniformly $o(N)$.
Thus this range, including its support-size sum, is exponentially small.
This finite-alphabet argument uses neither the general equality
classification nor the coarse ordinary subgroup-count estimate from
Theorem~\ref{thm:small-support}.

If $C>\sqrt N/4$ and $T>a/140$, then $C<284T$, so
$T>\sqrt N/1136$. The negative term of the second envelope therefore
has order at least $N^{3/2}$ and absorbs its error. The first envelope
is exponentially small uniformly for $a\geq1$. After summing the original
multiplicity weights, at most $(N+1)^2$ parameter pairs $(a,T)$ remain.
Their sum proves the theorem.
\end{proof}

The same weighted estimates give a direct extension to one natural
$S_3$ orbit, with critical orbits and the original thirteen noncritical
colours on the remaining points. Write $\mathbf p$ for the critical
profile of rank $R$, keep the noncritical profile $\mathbf q$ fixed,
and put $C=2a+4T>0$ and $N=R+C+1$. The physical degree is $2N+1$.
The exterior group is a $2$-group. If $(g,d)$ lifts a three-cycle and
$d^{2^k}=1$, its $4^k$th power is $(g,1)$. Thus every subgroup full
on the $S_3$ coordinate contains $A_3$ on that coordinate alone.
The sign quotient, followed by its inverse full preimage, identifies
its full models bijectively with those of
$(\mathbf p+\mathbf e_{C_2},\mathbf q)$. Every original critical and
noncritical coordinate remains full; no full projection onto either
entire class product is required.

The physical weight is still the original one. If
$\omega(\mathbf p,\mathbf q)$ denotes the reciprocal product of the
original normalizers and occurrence factorials, then
\[
 \frac{\omega(\mathbf p,\mathbf q)}6
 =\frac{p_{C_2}+1}{3}\,
       \omega(\mathbf p+\mathbf e_{C_2},\mathbf q).
\]
The factor includes both the marker normalizer six and the new C2
occurrence factorial. The map $\mathbf p\mapsto\mathbf p+\mathbf e_{C_2}$
is injective into the critical profiles of rank $R+1$, so positivity
bounds the original weighted sum by $(R+1)/3$ times the complete
shifted sum, with $\mathbf q$ unchanged. Its half-degree is exactly
$(R+1)+C=N$. Moreover,
$L_{2N+1}\geq(2N+1)!c_NG_N$, retaining the nonnegative additional
term of the odd coefficient. The preceding weighted estimates therefore
give a polynomial times $2^{-\Omega(N)}$ for each such bin.
There are at most $(N+1)^2$ pairs $(a,T)$ with $C+1\leq N$, and the
weight correction is at most $(N+1)/3$. Their cubic polynomial cost
is absorbed by the exponential decay. This bounds the actual union
of physical subgroups, allowing overlapping profile presentations.
It is a specialization to this finite alphabet and leaves the general
complete-family marker transfer of Section~\ref{sec:marker-complete}
unchanged.

Combining \eqref{eq:proved-mixture-envelopes} with
Propositions~\ref{prop:transport}--\ref{prop:absorb-decorations}
proves the same negligible-mass conclusion for every fixed finite menu
of simultaneous equal-degree quotient replacements retaining a positive
fraction of the original noncritical support. The original small-support
range is removed before paying its decoration fibre.
